\section{A common semantic boundary for the integrated design}
\label{sec:semantic-boundaries}

The nine expert responses, identified as N1--N9 in the integration ledger,
agree on the most important correction to the original plan: certification,
search coverage and useful execution are different obligations. This section
states their shared foundation once. Later sections retain the distinct
constructions, counterexamples and unresolved questions that depend on it.
The arguments are mathematical design arguments, not new Lean formalizations.

\subsection{Freeze the problem before proposing a solution}

Let the original query be
\[
 Q=(E,\Gamma,T,C,A,\mathcal L,\mathcal G,B,\Pi),
\]
where $E$ is the imported environment, $\Gamma$ the local telescope and its
instances, $T$ the target, $C:T\to\Prop$ the contract, $A$ the permitted axiom
policy, $\mathcal L$ the providers, and $(\mathcal G,B)$ a declared grammar
and its bounds. The frozen policy $\Pi$ identifies the ranking objective,
logical resource accounting, transparency/equality settings, search and
proposal admission policy, and their versions. A policy change is recorded
as a new experiment or an explicit epoch transition; it must not silently
reuse a certificate or cache entry whose assumptions changed.
An absent behavioral contract means $C(p)=\mathsf{True}$.
Closing a local answer abstracts the original telescope; it does not replace
its dependencies with similarly named variables. Original universe parameters
remain rigid. A term at $T[\ell:=0]$ answers a specialized query, not an
arbitrary original universe-polymorphic query. Any permitted generalization of
new inference variables must preserve the jointly checked program/proof types.

\begin{definition}[Completion set]\label{def:completion-set}
A partial state $S$ contains the typed program root, expression and universe
assignments, delayed constraints, local contexts and instances, recursion
capabilities, residual observations and proof obligations. Write
$\Comp_{\mathcal G,B}(S)$ for its compatible closing substitutions under the
declared grammar and bounds. A completion is a native well-scoped substitution,
not a textual replacement of question marks.
\end{definition}

Finiteness needs actual bounds on literals, universe expressions, auxiliary
declarations, recursion schemes and proof plans as well as term depth. A
finite program grammar alone does not turn a bounded tactic portfolio into a
decision procedure for arbitrary Lean propositions. Conversely, a finite
example conjunction can itself have a fully checked proof: its finite scope
does not make the proof less valid, but it does not specify an unstated
universal intention.

\subsection{Six guarantees and their separate evidence}

\begin{description}
\item[Positive certification.]
The closed abstractions of $p:T$ and $\pi:C(p)$ check in $E$, and the transitive
axiom dependencies of the program, original type, proof and original contract
are permitted by $A$. A kernel-checked theorem may use allowed axioms; it need
not be axiom-free.
\item[Conservative pruning.]
A refuted partial branch has no permitted satisfying completion. This
requires a separate argument about the refuter. Checking all surviving
answers cannot reveal a solution discarded before the gate.
\item[Relative coverage.]
Every member of a named candidate grammar is considered or safely excluded
by a fair exhaustive run. Certification completeness additionally needs a
decisive verifier for the advertised fragment, or a separate finite proof
grammar. Top-$k$ retrieval, queue eviction and a deadline are budget choices,
not exclusion proofs.
\item[Ranking optimality.]
No eligible candidate improves one fixed objective. First success, best among
visited terms, immutable program cost and a learned dispatch score are not
interchangeable. Outstanding alternatives need valid bounds before an
optimality claim is made.
\item[Replay.]
Proof replay, logical search replay and reproduction of model inference have
different prerequisites. A fixed logical prefix can be replayed; equal wall
time on different machines does not imply an equal completed prefix.
\item[Executable publication.]
Declaration transport, compilation and advertised execution are independently
checked. A well-typed recursor expression does not by itself establish that
the supported publication path compiles or executes it.
\end{description}

These are six obligations rather than six promised completed features.
The responses group them differently: most distinguish the first five,
N6 groups its principal guarantees more coarsely, and N5 and several
realization discussions make the sixth explicit. No disagreement is resolved by silently
folding execution into kernel acceptance.

\subsection{Necessary observations, alternatives and refutation}

For a partial root $p_S$, a sound hard rejection requires
\begin{equation}\label{eq:conservative-pruning}
 \forall\theta\in\Comp_{\mathcal G,B}(S),\quad \neg C(p_S\theta).
\end{equation}
An inferred observation or angelic constraint $D$ is safe to add to every
remaining alternative only if
\begin{equation}\label{eq:necessary-observation}
 \forall\theta\in\Comp_{\mathcal G,B}(S),\quad
 C(p_S\theta)\Longrightarrow D(\theta).
\end{equation}
If $D$ has no compatible realization, these two statements imply safe
refutation: a satisfying completion would satisfy $D$, contradicting its
impossibility. An observation need not hold for completions that violate the
contract. This is why an example can constrain a recursive call conditionally
on correctness of the whole synthesized function.

A sufficient guess has a different role. One guessed invariant, recursive
result or helper can define a search alternative; failure of that alternative
does not reject all other guesses. Repeated equal recursive calls must remain
correlated. Fuel exhaustion, an unsupported evaluator case, a failed tactic
or an unsuccessful unification attempt remains inconclusive unless a stronger
decisive fragment supplies the missing exclusion argument. Differential
testing can falsify a refuter's correctness, but passing tests cannot prove
Equation~\ref{eq:conservative-pruning} for every open completion.

\Needspace{27\baselineskip}
\subsection{Results must identify what was established}

\begin{longtable}{@{}L{32mm}L{119mm}@{}}
\caption{A result algebra with explicit evidence boundaries.}\label{tab:result-evidence}\\
\toprule Result & Required meaning\\\midrule\endfirsthead
\toprule Result & Required meaning\\\midrule\endhead
Verified & A term of the frozen target and a checked proof of its exact requested contract under its policy.\\
Observed / unrefuted & Recorded observations for a typed candidate; required certification is still missing.\\
Candidate refuted & A particular candidate or partial branch violates a checked necessary condition.\\
Impossible & A checked theorem excludes solutions of the original query under the advertised policy.\\
Bound completed & The named finite program/proof search was covered; unresolved verification obligations are recorded separately.\\
Interrupted / unsupported & Work stopped or a feature is unavailable; neither states that the target is uninhabited.\\
Publication status & Separate export, compilation and execution results for the accepted artifact.\\
\bottomrule
\end{longtable}

Refuting every visited candidate is not proving impossibility. Exhausting
syntax while some verification attempts remain unresolved is not semantic
exhaustion. A finite Kripke countermodel for a reified propositional calculus
establishes nonvalidity in that calculus; interpreting it as a Lean negation
of an arbitrary contextual target needs a further bridge. These distinctions
belong in receipts and diagnostics and need not become user-facing engine
settings.

\subsection{Services with deliberately different authority}

\begin{lstlisting}
ProposalService.propose : DecisionPoint -> SearchView -> Alternatives
ObservationService.check : PartialState -> ValuesOrBlockersOrRefutation
ProofService.tryProve : TypedObligation -> Budget -> ProofOrUnknown
AcceptanceGate.check : FrozenQuery -> Candidate -> Proof -> GateReply
PublicationService.realize : Accepted -> ExportCompileExecuteReceipt
\end{lstlisting}
This is interface pseudocode. A proposal changes the order or adds an
alternative; it carries no authority to remove every other alternative.
A refutation carries checked evidence or invokes a component with established
conservative semantics. The gate reconstructs $C(p)$ from the original query,
rather than trusting an unrelated supplied proof type. All services use the
same dependency information and retain their active resource checks.

The source packages inspected selected code at
\code{784664ff34e819422510a4d470ab07fe48bb736b}; their API observations are
historical. They correctly distinguish a gate accepting a proof/type pair
from an end-to-end caller that constructs the requested contract. Their
recommendation to freeze the query at the gate is not evidence that such a
caller accepted an unrelated proof. Likewise, N5's observations about negative
certificate retention, universe specialization and intermediate partial
reduction motivate explicit invariants; they must not be presented as
unrechecked defects in a later implementation. This documentation integration
does not change the implementation or claim a new acceptance run.

State ownership must be explicit across these services: immutable reusable
content, branch-local assignments/views, and monotone work/diagnostic counters
are different classes. An \code{IO.Ref} is not automatically a globally valid
cache. The joint dependency and scheduling obligations are developed in
Sections~\ref{sec:search} and~\ref{sec:evaluation}; proof and realization
obligations remain attached to the same original query throughout.
